% Provenance: paper_P3_group_size (contrast density, Wu pairing, bridge-to-DPO
% remnant), paper_P7_zvf_controller (controller, design rules), E-R2b artifacts
% (er2b_*.json), group_size.tex + group_size_reconcile.
\chapter{Group Size as a Schedule}
\label{ch:groupsize}

\section{Static sweeps are inconclusive}
A single-seed sweep over $G\in\{2,4,8,16,32\}$ on Qwen3-8B/GSM8K is
non-monotone in fixed-step last-10 reward ($G{=}4{:}0.52$, $G{=}32{:}0.44$,
$G{=}16{:}0.38$, $G{=}2{:}0.38$, $G{=}8{:}0.34$), and a token-budget-normalised
reanalysis shifts the apparent optimum rightward with total tokens (toward
$G\approx32$ at $T\ge16$M tokens). Neither view supports a universal winner;
the honest summary is that fixed-$G$ comparisons are confounded by what the
budget is held constant \emph{in}.

\section{The matched-budget experiment (Claim 2)}
The decisive design holds the \emph{rollout budget} fixed at 2{,}560 rollouts
per arm: $G{=}2\times160$ optimiser steps versus $G{=}16\times20$ steps
(batch 8, 512-token completions, LoRA rank 4, seeds 123 and 456):
\begin{center}
\begin{tabular}{@{}lcc@{}}
\toprule
 & $G{=}2\times160$ & $G{=}16\times20$ \\
\midrule
final train reward & $\approx 0.9$--$1.0$ (pool mastered) & $\approx 0.3$--$0.5$ (mid-learning) \\
late-run \ZVF{} & $0.75$--$1.0$ (all-correct wall) & $0.00$--$0.25$ \\
late-run gradient signal & effectively zero & sustained \\
\bottomrule
\end{tabular}
\end{center}
By the reward axis alone the $G{=}2$ arms win outright --- which is exactly
the trap: their final steps are spent on zero-gradient groups.
Small $G$ converts the budget into more optimiser steps early --- consistent
with the per-rollout accounting of Chapter~5 --- and then exhausts its own
signal as accuracy rises: the $p\to1$ wall of $h_G(p)$. Large $G$ pays for
contrast it does not need early and retains signal late. \textbf{Group size
controls which end of training starves; it is a schedule variable.}

\section{Relation to the contrastive reading of GRPO}
The $G{=}2$ versus $G{=}16$ pairing is exactly the pairing behind the
``GRPO is secretly contrastive/DPO-like'' equivalence claims. Our direct test
of the generalised claim ($G{=}4\sim G{=}32$) shows retention of the large-$G$
arm's accuracy falling from $97.6\%$ at $T{=}1$M tokens to $72.7\%$ at
$T{=}64$M: the equivalence holds only in the under-trained regime. The
matched-budget panel adds the trajectory view: at a fixed rollout budget the
small-$G$ arm can even \emph{lead} on train reward, but the equivalence
claims say nothing about signal availability --- and that is where the
small-$G$ lead dies. The fuller
bridge-to-DPO treatment from the group-size working paper remains
thesis-external future work.

\begin{figure}[t]
\centering
\begin{tikzpicture}[node distance=5mm and 7mm,
  stage/.style={rectangle,rounded corners,draw,align=center,minimum height=9mm,
                font=\scriptsize,inner sep=4pt},
  lab/.style={font=\scriptsize,align=center}]
\node[stage,fill=blue!10]  (early) {early training\\$p$ low--mid};
\node[stage,fill=blue!20,right=of early]  (mid) {mid training\\$p$ rising};
\node[stage,fill=blue!30,right=of mid]    (late) {endgame\\$p\to1$};
\node[lab,below=2mm of early] {small $G$ efficient:\\most groups mixed,\\
  more optimiser steps\\per rollout};
\node[lab,below=2mm of mid] {monitor the pair\\(reward, \ZVF{})\\with Wilson CI};
\node[lab,below=2mm of late] {small $G$ starves\\(all-correct wall);\\
  large $G$ keeps\\signal alive};
\draw[-{Stealth},thick] (early) -- (mid);
\draw[-{Stealth},thick] (mid) -- (late);
\draw[decorate,decoration={brace,amplitude=4pt},thick]
  ([yshift=3mm]early.north west) -- ([yshift=3mm]late.north east)
  node[midway,above=5pt,font=\scriptsize]
  {escalate $G$ when the T2 budget $N(\underline{\ZVF}_t)$ exceeds the step's
   rollout budget};
\end{tikzpicture}
\caption{Group size as a schedule. Per-rollout efficiency favours small $G$
early (Chapter~5); the all-correct wall punishes it late (Claim~2). The
escalation trigger comes from the T2 reliability budget, not from a
data-dependent ``optimal $G$'' --- which Chapter~5 shows does not exist under
per-rollout accounting.}
\label{fig:schedule}
\end{figure}

\section{Controller design (evidence-gapped, stated as such)}
A ZVF-triggered group-size controller follows from Chapters~4--5: measure
$\ZVF_t$ with its Wilson interval (T1); when the reliability budget
$N(\underline{\ZVF}_t)$ exceeds the step's remaining rollout budget, the step
fails to produce any informative group with probability $>\delta$ --- escalate
$G$ rather than gamble the budget (T2). Design rules from the controller
working paper carry over: gate on magnitude-aware statistics, escalate
multiplicatively in small steps, trigger on spikes rather than levels, never
compare controllers across stacks, and treat bought signal ($\ZVF{=}0$ via
dynamic sampling at ${\sim}45\%$ extra rollouts) as a purchase to be priced.
An offline bandit pilot recovers $97\%$ of a static oracle's informative
groups per 1{,}000 rollouts and beats uniform allocation by $28\%$.
\textbf{No efficacy claim is made} (non-claim 3): the decisive test --- a
pre-registered, compute-matched comparison against static $G{=}16$ and a
reward-only schedule, at $\ge3$ seeds on held-out metrics --- is future work,
and Chapter~5's T3 result removes any pretence that the controller could
derive a data-dependent target $G$ from the current objective.
